\section{Related Work}
\label{sec:related}

\para{Activation steering.}
The idea of controlling model behavior by intervening on internal representations has developed rapidly. \citet{turner2023steering} introduced Activation Addition (ActAdd), computing steering vectors from contrastive prompt pairs and adding them to residual stream activations during inference. \citet{rimsky2023steering} scaled this to multiple behavioral dimensions on Llama~2, showing that averaged contrastive vectors are effective across many behaviors and robust to system prompts. \citet{arditi2024refusal} demonstrated that refusal in chat models is mediated by a single direction, removable by projection. \citet{li2023inference} applied inference-time intervention to elicit truthful answers by shifting activations along truthful directions in attention heads. These methods share a common requirement: the user must articulate the target behavior well enough to construct contrastive examples or prompt pairs. We focus on the case where this articulation breaks down.

\para{Reliability of steering.}
\citet{braun2025unreliability} conducted a systematic study of steering vector reliability, finding that different prompt types produce directionally different vectors (low cosine similarity) and that steering is unreliable when the target behavior lacks a coherent directional representation across prompts. This directly motivates our work: if a feature cannot be described consistently, prompting-based steering will be inconsistent. We build on this finding by showing that learned embeddings can bypass the describability bottleneck.

\para{Soft prompt tuning and concept tokens.}
\citet{lester2021power} showed that learning soft prompts---continuous token embeddings prepended to the input---can match full finetuning at scale. \citet{li2021prefix} extended this to prefix-tuning, optimizing continuous vectors at every layer. Most relevant to our work, \citet{sastre2026concept} introduced \concepttokens: a single new token whose embedding is learned from natural language definitions, which can then steer behavior directionally. They showed that concept tokens outperform explicit naming for recasting (62\% vs.\ 17\%), a behavior that is difficult to describe precisely. Unlike \citet{sastre2026concept}, who train from natural language definitions, we train concept tokens to reference specific residual stream directions identified via contrastive activation analysis. This lets us study the relationship between a direction's describability and the benefit of using learned embeddings to access it. \citet{ackerman2024representation} proposed representation tuning, which permanently finetunes steering vectors into model weights using a dual loss function. Our approach is complementary: we learn a portable embedding rather than modifying the model.

\para{Linear representations in language models.}
The linear representation hypothesis \cite{park2024linear} posits that high-level concepts are encoded as linear directions in activation space. \citet{zou2023representation} formalized representation engineering as a framework for reading and controlling these directions. \citet{elhage2022superposition} showed that neural networks can represent more features than they have dimensions via superposition, implying that some directions may be difficult to isolate. \citet{shafran2026directions} argued that single directions are insufficient for many concepts, proposing region-based alternatives. Our work takes the linear representation hypothesis as a starting point and asks: given that a feature is linearly represented, does its accessibility via natural language affect how well it can be manipulated?
